\section{Coulomb's Law}
\subsection{Background}
Coulomb was a scientist who studied electricity in the early 1800s. He wanted to determine which factors affect 
\underline{the electrostatic force between two charged objects}. \\

Coulomb based his experiment on \underline{Cavendish's} experiment.\\

To perform the experiment, Coulomb needed to electrically charge each of the pith balls and 
know the \underline{magnitude of the charge} on each ball. His solution for this was to find the 
\underline{relative} magnitude of the charge on each pith ball. 

\subsection{Formula}
By measuring the force, the separation distance between the charged objects, and the relative 
charge on the pith balls, Coulomb was able to find the following relationships:
\begin{align}
    F_{E} &\propto \tfrac{1}{R^2} \nonumber \\
    F_{E} &\propto q_{1}q_{2} \nonumber
\end{align}

We can bring these proportionalities together:
\begin{equation*}
    \left| F_E \right| = \frac{k \left |q_1 \right |  \left |q_2 \right |}{R^2}
\end{equation*}
\begin{center}
    $F_E$ is the magnitude of the electrical force in between \underline{two point charges}\\
    $q_a$ and $q_b$ are the absolute values of the charge on each object (in C)\\
    $R$ is the separation distance between the objects (in m)\\
    $k$ is Coulomb's law constant of proportionality ($k = 8.99 \times 10^9 \tfrac{Nm^2}{C^2}$)
\end{center}

\begin{remark}
    When using equations for electrical forces, \underline{do not substitute the sign of the charge}. Determine the direction of the force \underline{conceptually}!
\end{remark}

\subsubsection{Point Charge}
If the charged objects had a reasonably large size, then the electrical force on each object would cause the charges to \underline{\textit{move inside}} each object. In a point charge, 
the charge is treated as unable to move around. 

\subsubsection{Compared to $F_g = \frac{Gm_1m_2}{R^2}$}
\textbf{Similar}
\begin{itemize}
    \item $F \propto \tfrac{x}{R^2}$
\end{itemize}
\textbf{Difference}
\begin{itemize}
    \item $k > G$
    \item $F_E$ can represent either attraction or repulsion.
\end{itemize}
